\documentclass[12pt, a4paper]{article}

\usepackage[UTF8]{ctex}
\usepackage{geometry}
\usepackage{graphicx}
\usepackage{float}
\usepackage{listings}
\usepackage{xcolor}
\usepackage{enumitem}
\usepackage{caption}

\geometry{left=2.5cm, right=2.5cm, top=3cm, bottom=3cm}

\lstset{
    language=Python,
    basicstyle=\small\ttfamily,
    keywordstyle=\color{blue},
    stringstyle=\color{red},
    commentstyle=\color{green!50!black},
    showstringspaces=false,
    breaklines=true,
    frame=single,
    numbers=left,
    numberstyle=\tiny\color{gray}
}

\setlength{\parindent}{0pt}

\begin{document}

\begin{center}
    {\LARGE \textbf{《计算机视觉》实验报告}} \\[1.5cm]
    {\large \textbf{姓名：周梓俊 \quad 学号：23120899}} \\[1cm]
    {\large \textbf{实验 3}} \\[1.5cm]
\end{center}

% ==================== 任务 1 ====================
\textbf{一、任务1}

\begin{enumerate}[label=\textbf{\alph*)}, leftmargin=*]

\item \textbf{核心代码：}

\begin{lstlisting}
# 转换为HSV空间
hsv = cv2.cvtColor(img, cv2.COLOR_BGR2HSV)

# 分离RGB通道和HSV通道
b, g, r = cv2.split(img)
h_ch, s_ch, v_ch = cv2.split(hsv)

# 对RGB三个通道画三维图
rows = np.arange(0, img.shape[0])
cols = np.arange(0, img.shape[1])
X, Y = np.meshgrid(cols, rows)

# 采样防止图太大卡死
step = max(1, max(img.shape[0], img.shape[1]) // 100)
X_s = X[::step, ::step]
Y_s = Y[::step, ::step]

ax1 = fig.add_subplot(131, projection='3d')
ax1.plot_surface(X_s, Y_s, r[::step, ::step], cmap='Reds')

ax2 = fig.add_subplot(132, projection='3d')
ax2.plot_surface(X_s, Y_s, g[::step, ::step], cmap='Greens')

ax3 = fig.add_subplot(133, projection='3d')
ax3.plot_surface(X_s, Y_s, b[::step, ::step], cmap='Blues')
\end{lstlisting}

\item \textbf{实验结果截图}

\begin{figure}[H]
    \centering
    \begin{minipage}{0.3\textwidth}
        \centering
        \includegraphics[width=\textwidth]{R_channel.png}
        \\[2pt]{\small R通道}
    \end{minipage}\hfill
    \begin{minipage}{0.3\textwidth}
        \centering
        \includegraphics[width=\textwidth]{G_channel.png}
        \\[2pt]{\small G通道}
    \end{minipage}\hfill
    \begin{minipage}{0.3\textwidth}
        \centering
        \includegraphics[width=\textwidth]{B_channel.png}
        \\[2pt]{\small B通道}
    \end{minipage}
\end{figure}

\begin{figure}[H]
    \centering
    \begin{minipage}{0.3\textwidth}
        \centering
        \includegraphics[width=\textwidth]{H_channel.png}
        \\[2pt]{\small H通道}
    \end{minipage}\hfill
    \begin{minipage}{0.3\textwidth}
        \centering
        \includegraphics[width=\textwidth]{S_channel.png}
        \\[2pt]{\small S通道}
    \end{minipage}\hfill
    \begin{minipage}{0.3\textwidth}
        \centering
        \includegraphics[width=\textwidth]{V_channel.png}
        \\[2pt]{\small V通道}
    \end{minipage}
\end{figure}

\begin{figure}[H]
    \centering
    \includegraphics[width=0.9\textwidth]{rgb_3d_surface.png}
    \\[4pt]{\small RGB三通道三维图}
\end{figure}

\item \textbf{实验小结}

在任务1中，我首先将图像从BGR颜色空间转换为HSV颜色空间，然后使用 cv2.split 分别分离了RGB三个通道和HSV三个通道并显示。接着利用 matplotlib 的 plot\_surface 函数对R、G、B三个通道分别绘制了三维表面图，其中x和y轴表示像素坐标，z轴表示该通道的像素值。为了防止图片过大导致渲染缓慢，对像素进行了等间距采样。

\end{enumerate}

\vspace{1cm}

% ==================== 任务 2 ====================
\textbf{二、任务2}

\begin{enumerate}[label=\textbf{\alph*)}, leftmargin=*]

\item \textbf{核心代码：}

\begin{lstlisting}
home_color = cv2.imread("home_color.png")
home_gray = cv2.cvtColor(home_color, cv2.COLOR_BGR2GRAY)

# 灰度直方图
gray_hist = cv2.calcHist([home_gray], [0], None, [256], [0, 256])

# 彩色直方图（B/G/R三通道分别计算）
colors = ('b', 'g', 'r')
for i, col in enumerate(colors):
    hist = cv2.calcHist([home_color], [i], None, [256], [0, 256])
    axes[1, 1].plot(hist, color=col)

# ROI直方图，ROI区域 x:50-100, y:100-200
mask = np.zeros(home_gray.shape[:2], dtype=np.uint8)
mask[100:200, 50:100] = 255
roi_img = cv2.bitwise_and(home_color, home_color, mask=mask)
roi_hist = cv2.calcHist([home_gray], [0], mask, [256], [0, 256])
\end{lstlisting}

\item \textbf{实验结果截图}

\begin{figure}[H]
    \centering
    \includegraphics[width=0.85\textwidth]{histogram_gray_color.png}
    \\[4pt]{\small 灰度直方图与彩色直方图}
\end{figure}

\begin{figure}[H]
    \centering
    \includegraphics[width=0.95\textwidth]{roi_histogram.png}
    \\[4pt]{\small ROI区域：原图、Mask、ROI提取结果、ROI直方图}
\end{figure}

\item \textbf{实验小结}

在任务2中，我使用 cv2.calcHist 函数分别计算了灰度图的直方图和彩色图的BGR三通道直方图，并将原图与对应直方图拼接在同一个窗口中显示。对于ROI区域，通过创建一个指定区域为白色的mask，利用 bitwise\_and 提取了感兴趣区域的图像，并计算了该区域的灰度直方图。最终将原图、mask图、ROI提取结果和ROI直方图放在同一个窗口内展示。

\end{enumerate}

\vspace{1cm}

% ==================== 任务 3 ====================
\textbf{三、任务3}

\begin{enumerate}[label=\textbf{\alph*)}, leftmargin=*]

\item \textbf{核心代码：}

\begin{lstlisting}
gray3 = cv2.cvtColor(img3, cv2.COLOR_BGR2GRAY)

# 直方图均衡化
equalized = cv2.equalizeHist(gray3)

# 计算均衡化前后的直方图
hist_before = cv2.calcHist([gray3], [0], None, [256], [0, 256])
hist_after = cv2.calcHist([equalized], [0], None, [256], [0, 256])
\end{lstlisting}

\item \textbf{实验结果截图}

\begin{figure}[H]
    \centering
    \includegraphics[width=0.85\textwidth]{histogram_equalization.png}
    \\[4pt]{\small 直方图均衡化前后对比}
\end{figure}

\item \textbf{实验小结}

在任务3中，我使用 cv2.equalizeHist 函数对灰度图像进行了直方图均衡化处理。均衡化前的直方图像素值分布集中在某些区间，均衡化后像素值分布更加均匀，图像的对比度得到了明显提升，暗部细节更加清晰。

\end{enumerate}

\end{document}
